% =============================================================================
% Practical SAT Solving - Uebungsblatt 2
% Dozent: Prof. Dr. Carsten Sinz
% Autor: Erlind Sejdiu
% =============================================================================
\documentclass[11pt, a4paper]{article}

% --- Sprache und Layout ---
\usepackage[ngerman]{babel}
\usepackage[utf8]{inputenc}
\usepackage[T1]{fontenc}
\usepackage{lmodern}
\usepackage{microtype}
\usepackage[left=2.0cm, right=2.0cm, top=1.8cm, bottom=1.8cm]{geometry}
\usepackage{parskip}

% --- Mathe und Tabellen ---
\usepackage{amsmath}
\usepackage{amssymb}
\usepackage{booktabs}
\usepackage{array}
\usepackage{xcolor}

% --- Hyperref ---
\usepackage{hyperref}
\hypersetup{
    colorlinks=true,
    linkcolor=black,
    citecolor=black,
    urlcolor=blue
}

% --- Kurzschreibweisen ---
\newcommand{\nx}[1]{\overline{x_{#1}}}
\newcommand{\n}[1]{\overline{#1}}

\begin{document}

% --- Header ---
\begin{flushleft}
    \textbf{Practical SAT Solving} -- WS 2026/27 \hfill \textbf{Erlind Sejdiu \& Johannes Ott}\\
    Prof. Dr. Carsten Sinz \hfill Hochschule Karlsruhe (HKA)\\
    \rule{\textwidth}{0.4pt}
\end{flushleft}

\begin{center}
    {\Large\bfseries Übungsblatt 2: Aufgaben 1 \& 2}
\end{center}
\vspace{0.2em}

% =============================================================================
\section*{Aufgabe 1: Resolutionsbeweis}
% =============================================================================

Gegeben ist die Formel
\begin{equation*}
    \underbrace{(x_3 \lor x_4 \lor \nx1 \lor x_5)}_{(1)} \land
    \underbrace{(\nx3 \lor x_4 \lor x_5)}_{(2)} \land
    \underbrace{(x_3 \lor \nx4 \lor \nx1)}_{(3)} \land
    \underbrace{(x_1 \lor x_2)}_{(4)} \land
    \underbrace{(x_1 \lor \nx2)}_{(5)} \land
    \underbrace{(\nx1 \lor \nx5)}_{(6)} \land
    \underbrace{(\nx3 \lor \nx4 \lor x_5)}_{(7)}
\end{equation*}

Durch wiederholte Resolution wird die leere Klausel $\square$ abgeleitet:

\begin{center}
\renewcommand{\arraystretch}{1.25}
\begin{tabular}{@{}r>{$}l<{$}ll@{}}
    \toprule
    Nr. & \text{Resolvente} & aus & Variable \\
    \midrule
    (8)  & (x_1)                    & (4), (5)   & $x_2$ \\
    (9)  & (\nx5)                   & (6), (8)   & $x_1$ \\
    (10) & (x_3 \lor \nx4)          & (3), (8)   & $x_1$ \\
    (11) & (\nx4 \lor x_5)          & (7), (10)  & $x_3$ \\
    (12) & (\nx3 \lor x_5)          & (2), (11)  & $x_4$ \\
    (13) & (x_3 \lor x_4 \lor x_5)  & (1), (8)   & $x_1$ \\
    (14) & (x_3 \lor x_5)           & (13), (11) & $x_4$ \\
    (15) & (x_5)                    & (12), (14) & $x_3$ \\
    (16) & \square                  & (9), (15)  & $x_5$ \\
    \bottomrule
\end{tabular}
\end{center}

Da die leere Klausel ableitbar ist, ist die Formel \textbf{unerfüllbar}.

\newpage
% =============================================================================
\section*{Aufgabe 2: Tseitin-Transformation}
% =============================================================================

Gegeben ist die Formel
\begin{equation*}
    F = \bigl(\nx1 \land \overline{(x_3 \Leftrightarrow x_2)}\bigr) \lor
        \bigl((x_3 \Rightarrow \nx4) \land (x_1 \Rightarrow (x_2 \land \nx3)) \land x_4\bigr)
\end{equation*}

\subsection*{Hilfsvariablen}

Für jede Teilformel wird eine neue Variable eingeführt:
\begin{center}
\renewcommand{\arraystretch}{1.2}
\begin{tabular}{@{}>{$}l<{$}@{\quad}>{$}l<{$}@{}}
    a \Leftrightarrow (x_3 \Leftrightarrow x_2) & d \Leftrightarrow (x_2 \land \nx3) \\
    b \Leftrightarrow (\nx1 \land \n{a})        & e \Leftrightarrow (x_1 \Rightarrow d) \\
    c \Leftrightarrow (x_3 \Rightarrow \nx4)    & g \Leftrightarrow (c \land e) \\
    h \Leftrightarrow (g \land x_4)             & f \Leftrightarrow (b \lor h)
\end{tabular}
\end{center}

\subsection*{Tseitin-KNF}

\begin{equation*}
    T(F) = f \land T_\text{def}(f) \land T_\text{def}(a) \land T_\text{def}(b) \land T_\text{def}(c)
    \land T_\text{def}(d) \land T_\text{def}(e) \land T_\text{def}(g) \land T_\text{def}(h)
\end{equation*}

\begin{center}
\renewcommand{\arraystretch}{1.25}
\begin{tabular}{@{}l>{$}l<{$}@{}}
    \toprule
    Definition & \text{Klauseln} \\
    \midrule
    Wurzel & (f) \\
    $f \Leftrightarrow (b \lor h)$ & (\n f \lor b \lor h) \land (f \lor \n b) \land (f \lor \n h) \\
    $a \Leftrightarrow (x_3 \Leftrightarrow x_2)$ &
        (\n a \lor \nx3 \lor x_2) \land (\n a \lor x_3 \lor \nx2) \land (a \lor x_3 \lor x_2) \land (a \lor \nx3 \lor \nx2) \\
    $b \Leftrightarrow (\nx1 \land \n a)$ & (\n b \lor \nx1) \land (\n b \lor \n a) \land (b \lor x_1 \lor a) \\
    $c \Leftrightarrow (x_3 \Rightarrow \nx4)$ & (\n c \lor \nx3 \lor \nx4) \land (c \lor x_3) \land (c \lor x_4) \\
    $d \Leftrightarrow (x_2 \land \nx3)$ & (\n d \lor x_2) \land (\n d \lor \nx3) \land (d \lor \nx2 \lor x_3) \\
    $e \Leftrightarrow (x_1 \Rightarrow d)$ & (\n e \lor \nx1 \lor d) \land (e \lor x_1) \land (e \lor \n d) \\
    $g \Leftrightarrow (c \land e)$ & (\n g \lor c) \land (\n g \lor e) \land (g \lor \n c \lor \n e) \\
    $h \Leftrightarrow (g \land x_4)$ & (\n h \lor g) \land (\n h \lor x_4) \land (h \lor \n g \lor \nx4) \\
    \bottomrule
\end{tabular}
\end{center}
\end{document}
